\section{Evaluación automática de lenguaje}
\textcolor{teal}{Revisión de técnicas de Automated Essay Scoring y Automatic Short Answer Grading. Revisión de los artículos académicos relacionados con el análisis y evaluación de lenguaje escrito.}

\textcolor{teal}{
    2.2.1 Automated Essay Scoring (AES)
    Qué explicar:
    \begin{itemize}
        \item Evaluación automática de textos largos
        \item Métodos tradicionales:
        \begin{itemize}
            \item features lingüísticas (longitud, vocabulario)
            \item modelos supervisados            
        \end{itemize}
    \end{itemize}
    Problemas:
    \begin{itemize}
        \item poca interpretabilidad
        \item dependencia de datasets anotados
        \item dificultad con contenido semántico
    \end{itemize}
    Transición a LLMs:
    \begin{itemize}
        \item LLMs permiten evaluar coherencia y generar feedback
    \end{itemize}
    Pero:
    \begin{itemize}
        \item riesgo de inconsistencia
        \item falta de control
    \end{itemize}
    2.2.2 Automatic Short Answer Grading (ASAG)
    Qué incluir:
    \begin{itemize}
        \item evaluación de respuestas cortas
        \item comparación con respuestas de referencia
    \end{itemize}
    Métodos:
    \begin{itemize}
        \item similitud semántica
        \item embeddings
        \item modelos supervisados
    \end{itemize}
    Problemas:
    \begin{itemize}
        \item múltiples respuestas correctas
        \item ambigüedad
    \end{itemize}
}